\documentclass{aa}

\usepackage{graphicx}
\usepackage{txfonts}
\usepackage{booktabs}
\usepackage{threeparttable}
\usepackage{siunitx}
\usepackage{amsmath}
\usepackage{hyperref}

\begin{document}

\title{The concentration-mass relation of cluster galaxies from strong lensing}

\author{M.~Meneghetti\inst{1}}

\institute{INAF -- Osservatorio di Astrofisica e Scienza dello Spazio di Bologna, via Gobetti 93/3, 40129 Bologna, Italy}

\date{Received --; accepted --}

\abstract
{We present a strong-lensing-based analysis of the concentration-mass relation of galaxy-scale halos in four massive clusters. The analysis is implemented with the scripts \texttt{fit\_bergamini\_bmo\_free\_jaffe\_grid.py} and \texttt{tutorial\_density\_profile\_fits.ipynb}.}
{Our goal is to map the PIEMD description of cluster members, commonly adopted in parametric lens models, onto a physically interpretable stellar+dark-matter decomposition and to infer the implied $c_{200}$--$M_{200}$ relation.}
{We model each galaxy total density as a singular PIEMD and fit it with the sum of a Jaffe stellar component and a truncated NFW-like dark-matter component in BMO form with outer slope matched to the PIEMD truncation behavior. We quantify fitting quality and derive $c_{200}$ and $M_{200}$ from the best-fit BMO parameters. We then calibrate empirical relations linking $(\sigma_0,r_{\rm cut})$ and $(c_{200},M_{200})$ for MACSJ0416, MACSJ1206, AS1063, and A2744.}
{The Jaffe+BMO decomposition reproduces PIEMD profiles with very small residuals (typical median $\mathrm{RMSE}_{\log\rho}\simeq 0.045$ dex and $R^2_{\log\rho}\simeq0.99998$ in all clusters). The inferred concentration levels are high, with median $c_{200}\sim 19$--$61$ over median masses $M_{200}\sim 0.9$--$2.9\times10^{11}\,M_\odot$, depending on cluster. We extend the empirical mapping with an explicit redshift term based on $\rho_{\rm crit}(z)$, which removes the large normalization mismatch found for high-$z$ external validation (El Gordo). Relative to LCDM benchmark relations (Duffy et al. and Dutton \& Macci\`o), observed concentrations remain larger by factors of $\sim2.3$--$9.6$ at fixed mass.}
{Within this modeling framework, galaxy-scale halos in cluster cores appear significantly over-concentrated compared with standard LCDM expectations. We also find tension between lensing-inferred stellar normalizations and photometric stellar-mass estimates (notably in MACSJ0416), suggesting either residual systematics in the PIEMD-based decomposition or in the adopted photometric mass calibration.}

\keywords{gravitational lensing: strong -- galaxies: clusters: general -- galaxies: halos -- dark matter -- methods: data analysis}

\section{Introduction}
Strong lensing by galaxy clusters provides one of the most direct constraints on the projected mass distribution in dense environments \citep{Bartelmann2001}. Modern parametric reconstructions represent the cluster mass as the superposition of a few large-scale halos and many galaxy-scale subhalos, allowing precise recovery of critical curves and multiple-image configurations \citep{Kneib1996,Jullo2007}.

A standard choice for the galaxy-scale component is the pseudo-isothermal elliptical mass distribution (PIEMD, or dPIE in the spherical limit), often with vanishing core radius for cluster members \citep{Jullo2007,Bergamini2019}. In this formalism, galaxy parameters are linked through luminosity-based scaling laws, reducing the dimensionality of the model and enabling efficient optimization against large sets of lensing constraints.

The main limitation is interpretability: PIEMD parameters are convenient for lens modeling, but they are not directly equivalent to physically motivated stellar and dark-matter decompositions. In particular, concentrations and virial masses are naturally defined for NFW-like halos \citep{NFW1997}, while stellar masses are more naturally tied to centrally concentrated baryonic profiles.

This work addresses that gap. Using the local analysis pipeline in \texttt{fit\_bergamini\_bmo\_free\_jaffe\_grid.py} and \texttt{tutorial\_density\_profile\_fits.ipynb}, we decompose PIEMD galaxy profiles into a Jaffe stellar component plus a truncated NFW-like dark-matter component (BMO form), then derive the implied $c_{200}$--$M_{200}$ relations for four clusters. The paper emphasizes methodology and reproducible quantitative outputs from these scripts.

\section{Methodology}

\subsection{Modeling of cluster galaxies}
Cluster strong-lensing models are built as a sum of:
\begin{enumerate}
\item large-scale cluster halos, and
\item a population of galaxy-scale halos associated with member galaxies.
\end{enumerate}

For the galaxy component, we adopt spherical PIEMD profiles with $r_{\rm core}\to 0$ (the choice used in the fitting workflow; \citealt{Jullo2007,Bergamini2019}):
\begin{equation}
\rho_{\rm PIEMD}(r)=\frac{\sigma_0^2}{2\pi G}\,\frac{r_{\rm cut}^2}{r^2(r^2+r_{\rm cut}^2)}.
\label{eq:piemd_singular}
\end{equation}
For completeness, the finite-core expression is
\begin{equation}
\rho_{\rm PIEMD}(r)=\frac{\sigma_0^2}{2\pi G}\frac{r_{\rm cut}+r_{\rm core}}{r_{\rm core}^2 r_{\rm cut}}\frac{1}{\left(1+r^2/r_{\rm core}^2\right)\left(1+r^2/r_{\rm cut}^2\right)}.
\label{eq:piemd_full}
\end{equation}

In parametric lensing reconstructions, $\sigma_0$ and $r_{\rm cut}$ are usually tied to galaxy luminosity $L$ through power-law scalings:
\begin{equation}
\sigma_0 = \sigma_{0,\star}\left(\frac{L}{L_\star}\right)^{\alpha},
\qquad
r_{\rm cut}=r_{{\rm cut},\star}\left(\frac{L}{L_\star}\right)^{\delta}.
\label{eq:lum_scalings}
\end{equation}
Eliminating $L$ gives
\begin{equation}
r_{\rm cut}=r_{{\rm cut},\star}\left(\frac{\sigma_0}{\sigma_{0,\star}}\right)^{\beta},
\qquad \beta=\delta/\alpha.
\label{eq:rcut_sigma}
\end{equation}
This is the relation used and measured in our workflow.

\subsection{Decomposing the PIEMD}
We model each PIEMD profile as
\begin{equation}
\rho_{\rm PIEMD}(r) \simeq \rho_{\rm Jaffe}(r;M_\star,r_J) + \rho_{\rm BMO}(r;\rho_s,r_s,r_t),
\label{eq:decomp}
\end{equation}
where the stellar component is \citep{Jaffe1983}
\begin{equation}
\rho_{\rm Jaffe}(r)=\frac{M_\star r_J}{4\pi r^2(r+r_J)^2},
\label{eq:jaffe}
\end{equation}
and the dark-matter component is a smoothly truncated NFW profile in BMO form \citep{Baltz2009}, based on the NFW family \citep{NFW1996,NFW1997}:
\begin{equation}
\rho_{\rm BMO}(r)=\frac{\rho_s}{x(1+x)^2}\left(\frac{\tau^2}{x^2+\tau^2}\right)^{1/2},
\qquad x\equiv r/r_s,\ \tau\equiv r_t/r_s.
\label{eq:bmo}
\end{equation}
With $n=1/2$ in Eq.~\ref{eq:bmo}, the asymptotic slope is close to $r^{-4}$, consistent with the large-radius PIEMD truncation behavior.

The fits are performed in log-density space over each galaxy radial grid, solving for $(r_s,r_t,\rho_s,M_\star,r_J)$. The quality metrics saved by the pipeline are $\mathrm{RMSE}_{\log\rho}$ and $R^2_{\log\rho}$. For the four cluster runs used here, the median values are:
\begin{equation}
\mathrm{median}(\mathrm{RMSE}_{\log\rho})\approx 0.045\ \mathrm{dex},
\qquad
\mathrm{median}(R^2_{\log\rho})\approx 0.99998,
\end{equation}
with small cluster-to-cluster scatter.

Figure~\ref{fig:mosaic} shows a 3\,$\times$\,3 subset of PIEMD profiles and best-fit Jaffe+BMO decompositions.

\begin{figure*}
\centering
\includegraphics[width=0.92\textwidth]{bergamini_bmo_free_jaffe_profiles_mosaic_3x3.png}
\caption{Examples of PIEMD profile decompositions into Jaffe (stellar) + BMO (dark matter), extracted from the fit products of \texttt{fit\_bergamini\_bmo\_free\_jaffe\_grid.py}. Each panel reports $\sigma_0$, $r_{\rm cut}$, best-fit $c_{200}$, $M_{200}$, and log-space fit RMSE.}
\label{fig:mosaic}
\end{figure*}

\subsection{Data}
We analyze reconstructions for MACSJ0416, MACSJ1206, AS1063, and A2744, following the Bergamini et al. model family \citep{Bergamini2019,Bergamini2023a,Bergamini2023b}. Redshifts and component counts are read from the local LensTool \texttt{.par} files used by the scripts. Multiple-image and source-family counts are taken from the corresponding published model papers.

Table~\ref{tab:data} summarizes the key properties used in this work.

\begin{table*}
\centering
\begin{threeparttable}
\caption{Cluster sample and model characteristics used in this analysis.}
\label{tab:data}
\begin{tabular}{lccccccc}
\toprule
Cluster & $z_{\rm lens}$ & Main data sets & $N_{\rm LS}$ & $N_{\rm gal}$ & $N_{\rm mult}$ & $N_{\rm src}$ & Reference \\
\midrule
MACSJ0416 & 0.3960 & HST imaging + VLT/MUSE spectroscopy & 10 & 212 & 237 & 88 & \citealt{Bergamini2023a} \\
MACSJ1206 & 0.4390 & HST imaging + VLT/MUSE spectroscopy & 7 & 258 & 82 & 27 & \citealt{Bergamini2019} \\
AS1063    & 0.3480 & HST imaging + VLT/MUSE spectroscopy & 5 & 222 & 55 & 20 & \citealt{Bergamini2019} \\
A2744     & 0.3072 & HST/JWST imaging + spectroscopy & 9 & 171 & 149 & 50 & \citealt{Bergamini2023b} \\
\bottomrule
\end{tabular}
\begin{tablenotes}
\item $N_{\rm LS}$ and $N_{\rm gal}$ are counted from the local \texttt{.par} realizations (components without and with \texttt{mag}, respectively). Multiple-image/source counts are literature values for the corresponding published model generation.
\end{tablenotes}
\end{threeparttable}
\end{table*}

\section{Results}

\subsection{Concentration-Mass relation}
From each best-fit BMO halo, we compute $R_{200}$, $M_{200}$, and $c_{200}=R_{200}/r_s$. The derived concentration-mass trends are cluster dependent and very tight in each run. The median values are:
\begin{itemize}
\item MACSJ0416: $\tilde{c}_{200}=22.75$, $\widetilde{M}_{200}=2.34\times10^{11}\,M_\odot$,
\item MACSJ1206: $\tilde{c}_{200}=61.37$, $\widetilde{M}_{200}=9.28\times10^{10}\,M_\odot$,
\item AS1063: $\tilde{c}_{200}=49.14$, $\widetilde{M}_{200}=1.18\times10^{11}\,M_\odot$,
\item A2744: $\tilde{c}_{200}=19.32$, $\widetilde{M}_{200}=2.90\times10^{11}\,M_\odot$.
\end{itemize}

The fitted log-slopes of the measured relations, $\log_{10}c_{200}=a+b\log_{10}(M_{200}/10^{12}M_\odot)$, are listed in Table~\ref{tab:rels} and span from $b\approx-0.23$ (AS1063) to $b\approx0$ (A2744).

\subsection{Parametric modeling of galaxy concentrations}
Figure~\ref{fig:sigrcut} shows the measured $\sigma_0$--$r_{\rm cut}$ relations. For each cluster we fit
\begin{equation}
r_{\rm cut}=r_{{\rm cut},\mathrm{ref}}\left(\frac{\sigma_0}{220\,\mathrm{km\,s^{-1}}}\right)^\beta.
\label{eq:sigrcut_fit}
\end{equation}

We then map $(\sigma_0,\beta,r_{{\rm cut},\mathrm{ref}},z)$ into $(M_{200},c_{200})$ using the two-stage empirical model implemented in \texttt{tutorial\_density\_profile\_fits.ipynb}. Defining
\begin{equation}
x\equiv\log_{10}\left(\frac{\sigma_0}{220}\right),\quad
u\equiv\log_{10}\left(\frac{r_{{\rm cut},\mathrm{ref}}}{10\,\mathrm{kpc}}\right),\quad
v\equiv\beta,
\end{equation}
we write
\begin{equation}
\log_{10}\left(\frac{M_{200}}{10^{12}M_\odot}\right)=a_M+b_Mx,
\quad
\log_{10}(c_{200})=a_c+b_cx,
\end{equation}
with
\begin{align}
a_M &= -0.5296 + 0.9376\,u + 0.0020\,v,\\
b_M &=  2.1215 - 0.1385\,u + 0.9502\,v,\\
a_c &=  1.6437 - 0.6875\,u - 0.00037\,v,\\
b_c &=  0.6894 + 0.00789\,u - 0.6941\,v.
\end{align}
To account for redshift dependence in the conversion between BMO parameters and $(M_{200},c_{200})$, we define
\begin{equation}
q(z)\equiv \log_{10}\left[\frac{\rho_{\rm crit}(z)}{\rho_{\rm crit}(z_{\rm ref})}\right],\qquad z_{\rm ref}=0.396,
\end{equation}
consistent with the expected role of evolving reference density and pseudo-evolution in halo-based quantities \citep{Diemer2013,DiemerKravtsov2015}.
and apply
\begin{align}
\log_{10}\left(\frac{M_{200}}{10^{12}M_\odot}\right) &= a_M+b_Mx+\gamma_M q(z),\\
\log_{10}(c_{200}) &= a_c+b_cx+\gamma_C q(z).
\end{align}
The fitted redshift coefficients in the notebook calibration are
\begin{equation}
\gamma_M=-0.0333,\qquad \gamma_C=-0.3444.
\end{equation}

The resulting predicted $c_{200}$--$M_{200}$ relations match the measured ones very well (global residuals $\mathrm{RMS}(\log M)\simeq 0.0034$ dex and $\mathrm{RMS}(\log c)\simeq 0.0013$ dex for the calibration set). Figure~\ref{fig:cm_emp} shows the agreement.

For external validation on previously unused clusters (PLCK-G287 and El Gordo), adding the redshift term reduces the median concentration residual from $0.0450$ dex (no-$z$ model) to $0.0066$ dex (redshift-aware model). In El Gordo specifically, the concentration mismatch decreases from $\mathrm{RMS}(\Delta\log c)\simeq 0.0866$ to $\simeq 0.0112$ dex.

\begin{figure}
\centering
\includegraphics[width=\hsize]{tutorial_sigma0_vs_rcut_bergamini_bmo_clusters.png}
\caption{$\sigma_0$--$r_{\rm cut}$ relations for the four clusters, from \texttt{tutorial\_density\_profile\_fits.ipynb}. Different $(r_{\rm cut,ref},\beta)$ pairs map into different $c_{200}$--$M_{200}$ families.}
\label{fig:sigrcut}
\end{figure}

\begin{figure}
\centering
\includegraphics[width=\hsize]{tutorial_c200_m200_bergamini_bmo_comparison.png}
\caption{Measured and empirical-predicted $c_{200}$--$M_{200}$ relations for each cluster (data-driven calibration only).}
\label{fig:cm_emp}
\end{figure}

\begin{table*}
\centering
\caption{Empirical relation parameters from the cluster-by-cluster fits.}
\label{tab:rels}
\begin{tabular}{lcccccc}
\toprule
Cluster & $N$ & $r_{\rm cut,ref}$ [kpc] & $\beta$ & $\mathrm{d}\log c/\mathrm{d}\log M$ & $\tilde{c}_{200}$ & $\widetilde{M}_{200}$ [$M_\odot$] \\
\midrule
MACSJ0416 & 201 & 38.40 & 2.000 & -0.1768 & 22.75 & $2.34\times10^{11}$ \\
MACSJ1206 & 151 & 8.20  & 2.281 & -0.2089 & 61.37 & $9.28\times10^{10}$ \\
AS1063    & 160 & 12.44 & 2.493 & -0.2312 & 49.14 & $1.18\times10^{11}$ \\
A2744     & 201 & 33.40 & 1.022 & -0.0049 & 19.32 & $2.90\times10^{11}$ \\
\bottomrule
\end{tabular}
\end{table*}

\section{Discussion}

\subsection{Overconcentration of cluster galaxies}
Figure~\ref{fig:cm_lcdm} compares the empirical cluster relations with standard LCDM concentration-mass prescriptions \citep{Duffy2008,DuttonMaccio2014}. At the median masses of our samples, observed concentrations exceed these expectations by factors of about 2.3--9.6 depending on the cluster.

The offset is largest in MACSJ1206 and AS1063, and smaller (but still significant) in MACSJ0416 and A2744. This cluster-to-cluster variation likely reflects a combination of genuine environmental differences and modeling systematics (e.g., parametric form rigidity, selection of subhalos strongly constrained by lensing, and central-region weighting).

\begin{figure}
\centering
\includegraphics[width=\hsize]{tutorial_c200_m200_bergamini_bmo_with_lcdm.png}
\caption{Empirical $c_{200}$--$M_{200}$ relations compared to LCDM benchmark trends (Duffy et al.; Dutton \& Macci\`o). The observed relations are systematically shifted to higher concentrations at fixed mass.}
\label{fig:cm_lcdm}
\end{figure}

\subsection{Inconsistency with stellar masses}
The Jaffe component provides a lensing-based estimate of the stellar normalization. For MACSJ0416, comparing the fitted stellar normalization with the photometric estimate (derived from the F160W-based calibration used in the pipeline, in the spirit of \citealt{Grillo2015}) yields a median ratio
\begin{equation}
\mathrm{median}\left(M_{\star,\mathrm{lens}}/M_{\star,\mathrm{phot}}\right)\simeq 1.96,
\end{equation}
with a broad distribution.

This mismatch indicates that either:
\begin{enumerate}
\item the PIEMD-based decomposition is forcing part of the central dark matter into the stellar term,
\item the assumed stellar profile family (Jaffe) is not fully capturing real inner structure, and/or
\item the photometric stellar-mass calibration is biased for these cluster members.
\end{enumerate}

Figure~\ref{fig:mstar} shows the global comparison from the BMO+Jaffe dual-fit outputs.

\begin{figure}
\centering
\includegraphics[width=\hsize]{tutorial_section9_1_mstar_comparison.png}
\caption{Stellar-mass comparison from Sect.~9.1 of \texttt{tutorial\_density\_profile\_fits.ipynb}: photometric estimate versus free-Jaffe fit (left), and distribution of $\log_{10}(M_{\star}^{\rm free}/M_{\star}^{\rm phot})$ (right).}
\label{fig:mstar}
\end{figure}

\section{Summary and conclusions}
We developed and applied a reproducible workflow to translate PIEMD cluster-member models into a physically interpretable Jaffe+BMO decomposition, and to derive the implied $c_{200}$--$M_{200}$ relation for galaxy-scale halos in cluster cores.

Main results:
\begin{itemize}
\item PIEMD profiles are fitted by Jaffe+BMO with excellent accuracy (median $\mathrm{RMSE}_{\log\rho}\approx 0.045$ dex, median $R^2_{\log\rho}\approx 0.99998$).
\item Cluster-specific $\sigma_0$--$r_{\rm cut}$ relations differ substantially, with $(r_{\rm cut,ref},\beta)$ varying from $(8.2,2.28)$ in MACSJ1206 to $(38.4,2.00)$ in MACSJ0416.
\item These differences map directly into distinct $c_{200}$--$M_{200}$ relations.
\item The empirical parameterization implemented in \texttt{tutorial\_density\_profile\_fits.ipynb} reproduces the measured relations with residuals at the $10^{-3}$ dex level for this dataset.
\item A redshift-aware correction based on $\rho_{\rm crit}(z)$ is required for robust transfer to higher-$z$ clusters; it strongly improves external validation (notably El Gordo).
\item In all clusters, inferred concentrations are significantly above standard LCDM expectations at fixed mass.
\item Stellar masses inferred from the lensing decomposition can disagree with photometric estimates, highlighting a key systematic that must be controlled in future analyses.
\end{itemize}

Future work should include a hierarchical model combining all clusters simultaneously, explicit treatment of covariance between scaling-law parameters, and tests with alternative stellar and dark-matter profile families.

\begin{thebibliography}{}
\bibitem[Baltz et al.(2009)]{Baltz2009}
Baltz, E. A., Marshall, P., \& Oguri, M. 2009, JCAP, 01, 015

\bibitem[Bartelmann \& Schneider(2001)]{Bartelmann2001}
Bartelmann, M., \& Schneider, P. 2001, Phys. Rep., 340, 291

\bibitem[Bergamini et al.(2019)]{Bergamini2019}
Bergamini, P., Rosati, P., Grillo, C., et al. 2019, A\&A, 631, A130

\bibitem[Bergamini et al.(2023a)]{Bergamini2023a}
Bergamini, P., Grillo, C., Rosati, P., et al. 2023a, A\&A, 674, A79

\bibitem[Bergamini et al.(2023b)]{Bergamini2023b}
Bergamini, P., Acebron, A., Grillo, C., et al. 2023b, ApJ, 952, 84

\bibitem[Diemer et al.(2013)]{Diemer2013}
Diemer, B., More, S., \& Kravtsov, A. V. 2013, ApJ, 766, 25

\bibitem[Diemer \& Kravtsov(2015)]{DiemerKravtsov2015}
Diemer, B., \& Kravtsov, A. V. 2015, ApJ, 799, 108

\bibitem[Duffy et al.(2008)]{Duffy2008}
Duffy, A. R., Schaye, J., Kay, S. T., \& Dalla Vecchia, C. 2008, MNRAS, 390, L64

\bibitem[Dutton \& Macci\`o(2014)]{DuttonMaccio2014}
Dutton, A. A., \& Macci\`o, A. V. 2014, MNRAS, 441, 3359

\bibitem[Grillo et al.(2015)]{Grillo2015}
Grillo, C., Karman, W., Suyu, S. H., et al. 2015, ApJ, 800, 38

\bibitem[Jaffe(1983)]{Jaffe1983}
Jaffe, W. 1983, MNRAS, 202, 995

\bibitem[Jullo et al.(2007)]{Jullo2007}
Jullo, E., Kneib, J.-P., Limousin, M., et al. 2007, New J. Phys., 9, 447

\bibitem[Kneib et al.(1996)]{Kneib1996}
Kneib, J.-P., Ellis, R. S., Smail, I., Couch, W. J., \& Sharples, R. M. 1996, ApJ, 471, 643

\bibitem[Navarro et al.(1996)]{NFW1996}
Navarro, J. F., Frenk, C. S., \& White, S. D. M. 1996, ApJ, 462, 563

\bibitem[Navarro et al.(1997)]{NFW1997}
Navarro, J. F., Frenk, C. S., \& White, S. D. M. 1997, ApJ, 490, 493
\end{thebibliography}

\end{document}
